\chapter*{前書き}

これは初等幾何学を系統的に述べようとするものではありません。
作者の私的なまとめです。

一階述語論理の幾何学の公理系としては、タルスキーによるものがありますが、公理系がいわば極めて圧縮されているので、論証的展開の初期段階ほど難しいということになる。
加えて、展開初期の頃はむしろ幾何学的直感が推論を邪魔するのである。

一方、こちらの公理系は、確かにタルスキーの公理系よりもだいぶ公理が多いですが、まだしも使いやすいと思います。
（それに、「公理」と銘打たれた部分に関しては単に実数の公理系相当物であり、そこは言ってみれば「常識」で、意味があるのは「公準」と題した部分なのである。）

しかしながら、一階述語論理の理論に宿命的なものとしての使いにくさも、やはりあります。
例えば、直線などが論理的対象ではないことです。
直線の関数をこの体系で考えることはできない。
直線を定める二点の関数を考えることはできる。
また、直線を擬似的に主語のように見せかけて命題を言い表すことはできなくもないが、述語に対するまごうかたなき項として扱うことはできない。

角の概念もまた問題です。
角 \(AOB\) と言ったとき、イメージとしては劣角と優角の二つあるわけです。
ところが最も標準的に考えた場合、角は頂点を共有する二つの半直線であり、しかしこれだけでは大小どちらの角であるのか定まらない。
その暗黙の指定は角 \(AOB\) を考えている人間の頭の中にのみある。
そこで、このような曖昧さを払拭すべく、最も即物的かつシンプルに角を定義すれば、三点のなす三角形の角としての角もしくは三点が一直線上にあるときの平角あるいは大きさゼロの角、となります。
そしてこの場合も、角そのものを直接定義することはできず、角の合同関係を通して間接的に定義されるのです。

このように、明らかに存在しているはずなのに一階述語論理では対象物として扱えないというもどかしさがありますが、その代わり、公理系はゲーデルの不完全性定理の意味で「完全」になります。
